\section{Robustness and Claim Audit}

The main overstatement risk is not a single parameter choice. It is the temptation to turn
a strong option backtest into a broad claim about option alpha. The paper controls that
risk by separating mathematical results, generated empirical diagnostics, and production
claims. A theorem is proved. A test-window number is tied to a generated artifact. A live
trading claim is explicitly out of scope.

\begin{table}[H]
\centering
\resizebox{\textwidth}{!}{\input{tables/claim_audit}}
\caption{Claim-audit summary. The table is generated by the empirical runner.}
\label{tab:claimaudit}
\end{table}

The evidence hierarchy is deliberately asymmetric. The closed-form tangency theorem and
the PSD covariance result are mathematical statements; they do not depend on the local
OPRA sample. The empirical Sharpe, regression alpha, regime results, and leave-one-out
tests are sample statements; they support the measurement framework but do not prove a
stable premium. The post-cost layer is still a conservative research simulation, not live
broker evidence. It includes spreads, explicit fees, slippage, borrow, capacity, simulated
margin funding, and assignment-risk flags, but it does not include live fills, live broker
margin preview, order routing, or broker reconciliation. The option purchase price is
always included in P\&L through signed premium weights.

\begin{result}[Residual risk is part of the model]
\label{res:residualrisk}
An option-only covariance model that omits residual option-return covariance is not a
conservative Markowitz analogue. It is a Greek-span assumption. The empirical model should
therefore include \(\Sigma_\varepsilon\) unless the residual term is shown to be
negligible out of sample.
\end{result}

The realized diagnostics support this result. Risk calibration is informative but not
production-tight. P\&L attribution leaves material residual terms. After the exact-expiry
PIT audit, the stock-control regression alpha is small and individual stock betas are not
small. The combined META/NVDA/TSLA leave-one-out test remains positive, but it still has
high equity beta. This blocks any claim that the current result is independent of the
large-cap equity regime. These are not defects to hide. They are the distinction between a
publishable measurement paper and an overfit strategy note.

The robustness checks imported from the broader cross-asset research design are therefore
the ones that fit the option-only question: random feasible portfolios, block-bootstrap
intervals, HAC regression inference, equity-control regressions, delta-matched equity
benchmarks, gross-versus-net cost diagnostics, regime splits, and leave-one-underlying-out
tests. The verifier is intentionally asymmetric: it downgrades claims when VIX settlement,
net returns, inference, or data lineage do not support the stronger wording.
